\documentclass[a4paper,14pt]{extarticle}

\usepackage[left=30mm, right=10mm, top=20mm, bottom=20mm]{geometry}

\usepackage[T2A]{fontenc}
\usepackage[utf8]{inputenc}
\usepackage[russian]{babel}
\usepackage{tempora}
\usepackage{type1cm}
\usepackage{indentfirst}

\usepackage{minted}
\usepackage{xcolor}

\usepackage{siunitx}
\usepackage{booktabs}
\usepackage{multirow}


\usepackage{csquotes}
\usepackage[
    parentracker=true,
    backend=biber,
    style=gost-numeric,
    sorting=nty
]{biblatex}

\usepackage[nodisplayskipstretch]{setspace}
\usepackage[none]{hyphenat}
\usepackage{fancyhdr}
\usepackage{titlesec}
\usepackage{tocloft}

\usepackage{amssymb}
\usepackage{amsmath}
\usepackage{tikz}
\usepackage{pgfplots}
\usepackage{graphicx}
\usepackage{caption}
\usepackage{subcaption}
\usepackage[hidelinks]{hyperref}
\usepackage[russian]{cleveref}
\usepackage{bookmark}

\pgfplotsset{compat=1.18}

\onehalfspacing
\sloppy

\pagestyle{fancy}
\fancyhf{}
\fancyfoot[R]{\thepage}
\renewcommand{\headrulewidth}{0pt}

\definecolor{codebg}{rgb}{0.12,0.12,0.14}
\definecolor{codeframe}{rgb}{0.40,0.40,0.45}
\definecolor{keywordcolor}{rgb}{0.45,0.65,1.0}
\definecolor{stringcolor}{rgb}{1.0,0.55,0.45}
\definecolor{commentcolor}{rgb}{0.55,0.75,0.60}
\definecolor{textcolor}{rgb}{0.90,0.90,0.90}

\renewcommand{\theFancyVerbLine}{
	\sffamily\small\color{black}\arabic{FancyVerbLine}
}

\setminted{
    bgcolor=codebg,
	fillcolor=codebg,
	framesep=1em,
	frame=single,
    rulecolor=\color{codeframe},
    fontsize=\normalsize,
    baselinestretch=1.1,
    linenos,
    numbersep=6pt,
    tabsize=4,
    breaklines,
    breakbytoken,
	breaksymbolleft={},
    autogobble,
    showspaces=false,
    showtabs=false,
}

\usemintedstyle{lightbulb}
\setmintedinline{
    bgcolor=codebg
}

\makeatletter
\setlength{\@fptop}{0pt}
\setlength{\@fpsep}{8pt}
\setlength{\@fpbot}{0pt plus 1fil}
\makeatother

\newcommand{\circledequals}{\mathbin{\tikz \node[draw,circle,inner sep=1pt] {$=$};}}

\begin{document}
    \begin{titlepage}
		\centering
		
		\includegraphics{images/FEFU-logo.jpg}\\
		МИНИСТЕРСТВО ОБРАЗОВАНИЯ И НАУКИ РОССИЙСКОЙ ФЕДЕРАЦИИ\\
		Федеральное государственное автономное образовательное учреждение высшего образования\\
		\textbf{«Дальневосточный федеральный университет»}

		\rule{\linewidth}{0.8mm}\\[0.3cm]
		
		\textbf{ИНСТИТУТ МАТЕМАТИКИ И КОМПЬЮТЕРНЫХ ТЕХНОЛОГИЙ}\\
		\textbf{Кафедра математического и компьютерного моделирования}\\
		\vspace*{0.9cm}
		
		Кулахсзян Сергей Грайрович\\
        \vspace*{0.3cm}

		ЛАБОРАТОРНАЯ РАБОТА № 5 \\
        ЗАДАЧА № 32 \\
		\vspace*{0.5cm}

		\textbf{Задачи на электрические цепи}\\
		\vspace*{0.5cm}

		Направление подготовки 02.03.01сцт Сквозные цифровые технологии,\\ 
        бакалавриатская программа «Математика и компьютерные науки»\\ 
        Очной формы обучения\\
		\vspace{1.5cm}
		
		\vfill
		\centering
		г. Владивосток\\
		2026
	\end{titlepage}

    \tableofcontents
    \thispagestyle{empty}

	\newpage
    \pagenumbering{arabic}
    \setcounter{page}{3}
    \section{Постановка задачи}
    Идеальный колебательный контур с индуктивностью 1 Гн и ёмкостью 100 пФ подключается к источнику переменного напряжения $U(t) = U_0 \sin \omega t$.

    При какой частоте источника $\omega$ колебания тока будут иметь неограниченно возрастающую со временем амплитуду?

    \section{Математическое моделирование}
    Пусть $I(t)$ -- сила тока (в амперах) в цепи, $q(t)$ -- заряд (в кулонах) конденсатора в момент времени $t$ (в секундах).
    Напряжение на катушке определяется по формуле:
    \begin{equation} \label{U_coil}
        U(t) = LI'(t),
    \end{equation}
    где $L$ -- индуктивность (в генри).
    Напряжение на конденсаторе определяется соотношением:
    \begin{equation} \label{U_capacitor}
        \displaystyle U(t) = \frac{q(t)}{C},
    \end{equation}
    где $C$ -- ёмкость (в фарадах). Причём $q'(t) = I(t)$.

    Активный элемент ЭДС генерирует переменное напряжение:
    \begin{equation} \label{U_EMF}
        U(t) = U_0 \sin(\omega t + \phi),
    \end{equation}
    где $U_0$ -- начальное напряжение, $\omega$ -- частота, $\phi$ -- начальная фаза.

    Второй закон Кирхгофа гласит, что в цепи, представляющей собой замкнутый контур с единственным источником напряжения, сумма падений напряжений на пассивных элементах равна напряжению, вырабатываемому источником ЭДС.
    Тогда используя (\ref{U_coil}), (\ref{U_capacitor}) и (\ref{U_EMF}), получаем, что:
    \begin{equation} \label{2_Kirhgof_law}
        \displaystyle U_0 \sin (\omega t + \phi) = LI'(t) + \frac{q(t)}{C}.
    \end{equation}

    Дифференцируя (\ref{2_Kirhgof_law}), получаем математическую модель силы тока в цепи:
    \begin{equation}
        LI''(t) + \frac{I(t)}{C} = U_0 \omega \cos(\omega t + \phi)
    \end{equation}

    При подстановке значений коэффициентов из задачи получаем линейное неоднородное уравнение с постоянными коэффициентами:
    \begin{equation} \label{model}
        I''(t) + 10^{10} I(t) = U_0 \omega \cos \omega t.
    \end{equation}
    При решении уравнения получаем значение силы тока, из которой выводится амплитуда.

    \section{Решение}
    Уравнение (\ref{model}) имеет следующее общее решение: \\
    $\displaystyle
    \lambda^2 + 10^{10} = 0 \\
    \lambda_{1, 2} = \pm 10^5 i \\
    I = c_1 \cos 10^5 t + c_2 \sin 10^5 t
    $ -- это общее решение однородного уравнения, представляющее собой колебания с амплитудой $\sqrt{c_1^2 + c_2^2}$.

    Если $\omega \neq 10^5$, то частное решение ищется в виде
    $I = A \cos \omega t + B \sin \omega t$. 
    Это также незатухающие колебания постоянной амплитуды $\sqrt{A^2 + B^2}$.
    Следовательно, данный случай не удовлетворяет условию задачи.

    Если $\omega = 10^5$, то можно подобрать числа $A$ и $B$, при которых частное решение запишется в виде:
    \begin{equation} \label{particular_solution}
        I = t(A \cos 10^5 t + B \sin 10^5 t).
    \end{equation}
    Выражение (\ref{particular_solution}) задаёт колебания, имеющие непрерывно возрастающую с течением времени амплитуду $t\sqrt{A^2 + B^2}$.
    Тогда сила тока в цепи:
    \begin{equation*}
        I = (c_1 + At) \cos 10^5 t + (c_2 + Bt) \sin 10^5 t.
    \end{equation*}

    \section{Заключение}
    Тогда получаем, что, исходя из решения, источник ЭДС должен генерировать напряжение частотой 100000 Гц, чтобы амплитуда непрерывно росла.
\end{document}